\section{Well-formedness}
\label{sec:well-formedness}

\PurePy excludes some Python behaviours statically and others only at runtime.
Well-formedness is the static part, a judgement on programs defined in this \seckind{}; the operational
semantics (\secref{operational-semantics}) is defined for well-formed programs. A program is well-formed
when it is well typed, so the rules of this \seckind{} give each expression a type, each statement a
static outcome recording what it assigns or returns, and each module a signature. The \seckind{} introduces
types and subtyping, then follows the syntax of \secref{syntax}, taking statements, expressions and
patterns, then modules and programs, and \secref{operational-semantics} gives each of these a dynamic
counterpart.

\specinput{spec/well-formedness/auxiliary}
\paperinput{paper/well-formedness/principles}

\subsection{Types}
\label{sec:types}

\begin{figure}
\flushleft \shadebox{$\Gamma \vdash q : \theta$}
\begin{mathpar}
\small
\inferrule*[lab={\ruleName{simple-var}}]
{
  \bind{x}{\tau} \in \Gamma
}
{
  \Gamma \vdash x : \tau
}
\and
\inferrule*[lab={\ruleName{simple-module}}]
{
  \bind{x}{\modLoaded{q}{\Gamma'}} \in \Gamma
}
{
  \Gamma \vdash x : \modLoaded{q}{\Gamma'}
}
\and
\inferrule*[lab={\ruleName{simple-class}}]
{
  \bind{c}{\Type{C}} \in \Gamma
}
{
  \Gamma \vdash c : \Type{C}
}
\and
\inferrule*[lab={\ruleName{qualified}}]
{
  \Gamma \vdash q : \modLoaded{q'}{\Gamma'} \\
  \Gamma' \vdash x : \theta
}
{
  \Gamma \vdash q.x : \theta
}
\end{mathpar}
\caption{Qualified name $q$ resolves to context entry $\theta$ under $\Gamma$}
\label{fig:resolve-name}
\end{figure}


\begin{figure}
\flushleft \shadebox{$\replaced{\Sigma; \Gamma}{\Gamma} \vdash \psi \synth \tau$}
\begin{mathpar}
\small
\inferrule*[lab={\ruleName{ty-name}}]
{ \bind{\nu}{\predefName} \in \Gamma }
{ \replaced{\Sigma; \Gamma}{\Gamma} \vdash \nu \synth \nu }
\and
\inferrule*[lab={\ruleName{ty-var}}]
{
  \bind{\alpha}{\TypeVar} \in \Gamma
}
{
  \replaced{\Sigma; \Gamma}{\Gamma} \vdash \alpha \synth \alpha
}
\and
\inferrule*[lab={\ruleName{ty-class}}]
{
  \Gamma \vdash q : \Type{C} \\
  \added{\Sigma; \Gamma \vdash \seq{\psi} \synth \seq{\tau}} \\
  \added{C \text{ has } |\seq{\tau}| \text{ type parameters in } \Sigma}
}
{
  \replaced{\Sigma; \Gamma}{\Gamma} \vdash \replaced{q[\seq{\psi}]}{q} \synth \replaced{C[\seq{\tau}]}{C}
}
\and
\inferrule*[lab={\ruleName{ty-alias}}]
{
  \added{\Gamma \vdash q : \TypeAlias{\seq{\alpha}}{\tau}} \\
  \added{\Sigma; \Gamma \vdash \seq{\psi} \synth \seq{\sigma}}
}
{
  \added{\Sigma; \Gamma \vdash q[\seq{\psi}] \synth \subst{\seq{\sigma}}{\seq{\alpha}}{\tau}}
}
\and
\inferrule*[lab={\ruleName{ty-literal}}]
{ \bind{\pyinline{Literal}}{\predefName} \in \Gamma }
{ \replaced{\Sigma; \Gamma}{\Gamma} \vdash \tyLiteral{\ell} \synth \tyLiteral{\ell} }
\and
\inferrule*[lab={\ruleName{ty-list}}]
{ \bind{\pyinline{list}}{\predefName} \in \Gamma \\ \replaced{\Sigma; \Gamma}{\Gamma} \vdash \psi \synth \tau }
{ \replaced{\Sigma; \Gamma}{\Gamma} \vdash \tyList{\psi} \synth \tyList{\tau} }
\and
\inferrule*[lab={\ruleName{ty-dict}}]
{ \bind{\pyinline{dict}}{\predefName} \in \Gamma \\ \bind{\pyinline{str}}{\predefName} \in \Gamma \\ \replaced{\Sigma; \Gamma}{\Gamma} \vdash \psi \synth \tau }
{ \replaced{\Sigma; \Gamma}{\Gamma} \vdash \tyDict{\psi} \synth \tyDict{\tau} }
\and
\inferrule*[lab={\ruleName{ty-tuple}}]
{ \bind{\pyinline{tuple}}{\predefName} \in \Gamma \\ \replaced{\Sigma; \Gamma}{\Gamma} \vdash \psi_i \synth \tau_i \quad (\forall i) }
{ \replaced{\Sigma; \Gamma}{\Gamma} \vdash \tyTuple{\seq{\psi}} \synth \tyTuple{\seq{\tau}} }
\and
\inferrule*[lab={\ruleName{ty-callable}}]
{ \bind{\pyinline{Callable}}{\predefName} \in \Gamma \\ \replaced{\Sigma; \Gamma}{\Gamma} \vdash \psi_i \synth \sigma_i \quad (\forall i) \\ \replaced{\Sigma; \Gamma}{\Gamma} \vdash \psi \synth \tau }
{ \replaced{\Sigma; \Gamma}{\Gamma} \vdash \tyCallable{\seq{\psi}}{\psi} \synth \tyCallable{\seq{\sigma}}{\tau} }
\and
\inferrule*[lab={\ruleName{ty-union}}]
{ \replaced{\Sigma; \Gamma}{\Gamma} \vdash \psi \synth \sigma \\ \replaced{\Sigma; \Gamma}{\Gamma} \vdash \psi' \synth \tau }
{ \replaced{\Sigma; \Gamma}{\Gamma} \vdash \tyUnion{\psi}{\psi'} \synth \tyUnion{\sigma}{\tau} }
\end{mathpar}
\caption{Type expression $\psi$ denotes type $\tau$ under $\Gamma$\added{ and $\Sigma$}}
\label{fig:resolve-type}
\end{figure}


The rules are intended as a refinement of mypy's: a well-typed
\PurePy program should be accepted by mypy, and \PurePy may reject programs mypy accepts but not the
reverse. Function signatures, assignments and dataclass fields are annotated with type expressions, and a
type expression denotes a type under a context (\figref{resolve-type}). Every definition carries an
annotation for its parameters and return type, so the type of a definition is given rather than
inferred from its body. A class name resolves to a
class (\figref{resolve-name}) and every other name must resolve to a predefined entry
(\figref{predefined-types}). The function $\baseType(\tau)$ discards the precision of a
literal type, taking $\tyLiteral{42}$ to \pyinline{int} and leaving every other type alone
(\secref{types-auxiliary}).

\speconly{\subsubsection{Subtyping}\label{sec:subtyping}}%
\paperonly{\paragraph{Subtyping.}}

Subtyping (\figref{subtyping}) has every type below itself, a literal type below its base type, and a class below its
ancestors, with \pyinline{object} above every type and \pyinline{Never} below every type. A union is above
each of its members and below every type that is above all of them. A tuple type is covariant, below another of the
same length whose components are above its own. Two types are \emph{equivalent}, $\sigma \tyEquiv \tau$, when each is a subtype of the other; both notations follow \citet{pierce2002}. Lists and
dictionaries are related only where their components are equivalent, since both are invariant in mypy
and covariance would accept programs mypy rejects. A callable type is contravariant in its parameters
and covariant in its result, following mypy. The types with a length, namely lists, dictionaries,
strings and tuples, are below \pyinline{Sized}. An \pyinline{int} is a \pyinline{float}, as PEP~484
requires, so \pyinline{x: float = 0} is well typed. A \pyinline{bool} is not an \pyinline{int}, although
Python treats it as one.

\begin{lemma}[Transitivity of subtyping]
\label{lem:subtyping-transitive}
If $\Sigma \vdash \sigma \subty \tau$ and $\Sigma \vdash \tau \subty \tau'$ then $\Sigma \vdash \sigma \subty \tau'$.
\end{lemma}

\noindent Subtyping is therefore a preorder: reflexive and transitive but not antisymmetric,
because distinct types such as $\tyUnion{\sigma}{\tau}$ and $\tyUnion{\tau}{\sigma}$ are
equivalent.

Two types are \emph{comparable}
when either is a subtype of the other. Two comparable types
\emph{join} to the larger type and two incomparable types to their union. Two comparable types \emph{meet} to the smaller type; a
meet of incomparable types distributes over a union, takes tuples componentwise and callables with
parameters joined and results combined using meet, and is \pyinline{Never} otherwise, since no value belongs to both types.
So joining a type with itself gives that type, joining a class with one of its ancestors gives the ancestor, and
joining \pyinline{int} with \pyinline{str} gives $\tyUnion{\tyInt}{\tyStr}$. A sequence of types is joined by
folding binary join over the sequence, with \pyinline{Never} as the unit (\secref{types-auxiliary}).

\begin{lemma}[Types form a lattice]
\label{lem:lattice}
Types ordered by subtyping in $\Sigma$, up to equivalence, form a bounded distributive lattice: $\joinTy{\Sigma}{\sigma}{\tau}$
is the least upper bound of $\sigma$ and $\tau$, $\meetTy{\Sigma}{\sigma}{\tau}$ is their greatest lower bound,
\pyinline{object} is the greatest type and \pyinline{Never} the least.
\end{lemma}

\begin{figure}
\flushleft \shadebox{$\Gamma \vdash \sigma \subty \tau$}
\begin{mathpar}
\small
\inferrule*[lab={\ruleName{subty-refl}}]
{
  \strut
}
{
  \Gamma \vdash \tau \subty \tau
}
\and
\inferrule*[lab={\ruleName{subty-class}}]
{
  \Gamma(c') \in \ancestors(\Gamma(c))
}
{
  \Gamma \vdash q.c \subty q.c'
}
\and
\inferrule*[lab={\ruleName{subty-never}}]
{
  \strut
}
{
  \Gamma \vdash \tyNever \subty \tau
}
\and
\inferrule*[lab={\ruleName{subty-object}}]
{
  \strut
}
{
  \Gamma \vdash \tau \subty \tyObject
}
\and
\inferrule*[lab={\ruleName{subty-int-float}}]
{
  \strut
}
{
  \Gamma \vdash \tyInt \subty \tyFloat
}
\and
\inferrule*[lab={\ruleName{subty-literal}}]
{
  \Gamma \vdash \widen(\ell) \subty \tau
}
{
  \Gamma \vdash \tyLiteral{\ell} \subty \tau
}
\and
\inferrule*[lab={\ruleName{subty-union-left}}]
{
  \Gamma \vdash \sigma \subty \tau
}
{
  \Gamma \vdash \sigma \subty \tyUnion{\tau}{\tau'}
}
\and
\inferrule*[lab={\ruleName{subty-union-right}}]
{
  \Gamma \vdash \sigma \subty \tau'
}
{
  \Gamma \vdash \sigma \subty \tyUnion{\tau}{\tau'}
}
\and
\inferrule*[lab={\ruleName{subty-union}}]
{
  \Gamma \vdash \tau \subty \sigma \\
  \Gamma \vdash \tau' \subty \sigma
}
{
  \Gamma \vdash \tyUnion{\tau}{\tau'} \subty \sigma
}
\end{mathpar}
\caption{Subtyping}
\label{fig:subtyping}
\end{figure}



\subsection{Expressions}
\label{sec:expression-typing}

Within an expression the typing judgement gives each subexpression a type, and
is bidirectional. \emph{Synthesis} $\Sigma; \Gamma \vdash e \synth \tau$ reads a type out of an expression (\figrefThree{typing-synth}{typing-synth-elim}{typing-synth-ctd}),
and \emph{checking} $\Sigma; \Gamma \vdash e \checks \tau$ tests an expression against a type supplied by its
surroundings (\figref{typing-check}). The division follows the usual recipe for bidirectional
systems~\cite{davies2000,dunfield2021}: introduction forms check, since a lambda's parameters cannot be
annotated and an empty list does not determine its element type, and elimination forms synthesise, since a call, a subscript or an attribute
reference starts from a subterm whose type is known. Type systems that follow the recipe strictly usually
rely on an annotation form $(e \mathbin{:} \tau)$, which synthesises by construction, so a checking form
written with an annotation can be used in any context where synthesis is required. Python has no such form,
allowing annotations only on parameters, returns and assignments, so \PurePy gives the introduction forms
synthesis rules of their own wherever their parts determine the type; Swift, Kotlin and Scala
take similar approaches for similar reasons~\cite{swift2025,kotlin2025,odersky2025}.

Variables, literals, attribute references, calls, constructor calls, operators and subscripts synthesise.
The \ruleName{var} rule applies whenever $\Gamma(x)$ is a type, so a variable that is not definitely
assigned has no rule. A
literal $\ell$ synthesises the literal type $\tyLiteral{\ell}$.

\begin{figure}
\flushleft \shadebox{$\Sigma; \Gamma \vdash e \synth \tau$}
\begin{mathpar}
\small
\inferrule*[lab={\ruleName{var}}]
{
  \bind{x}{\tau} \in \Gamma
}
{
  \Sigma; \Gamma \vdash x \synth \tau
}
\and
\inferrule*[lab={\ruleName{literal}}]
{
  \strut
}
{
  \Sigma; \Gamma \vdash \ell \synth \tyLiteral{\ell}
}
\and
\inferrule*[lab={\ruleName{attr-module}}]
{
  \Gamma \vdash \exAttr{e}{x} : \tau
}
{
  \Sigma; \Gamma \vdash \exAttr{e}{x} \synth \tau
}
\and
\inferrule*[lab={\ruleName{attr-object}}]
{
  \Sigma; \Gamma \vdash e \synth \replaced{C[\seq{\tau}]}{C} \\
  \bind{x}{\sigma} \in \replaced{\tyInst{\fields(\Sigma, C)}{\seq{\tau}}}{\fields(\Sigma, C)}
}
{
  \Sigma; \Gamma \vdash \exAttr{e}{x} \synth \sigma
}
\and
\inferrule*[lab={\ruleName{call}}]
{
  \Sigma; \Gamma \vdash e \synth \tyCallable{\seq{\tau}}{\tau} \\
  \Sigma; \Gamma \vdash \seq{e} \checks \seq{\tau}
}
{
  \Sigma; \Gamma \vdash \exCall{e}{\seq{e}} \synth \tau
}
\and
\inferrule*[lab={\ruleName{syn-call-lambda}}]
{
  \Sigma; \Gamma \vdash \seq{e} \synth \seq{\tau} \\
  \Sigma; \Gamma \override \set{\bind{\seq{x}}{\seq{\tau}}} \vdash e \synth \tau
}
{
  \Sigma; \Gamma \vdash \exCall{(\exLambda{\seq{x}}{e})}{\seq{e}} \synth \tau
}
\and
\inferrule*[lab={\ruleName{constr}}]
{
  \Gamma \vdash q : \Type{C} \\
  \fields(\Sigma, C) = \bind{\seq{x}}{\seq{\tau}} \\
  \fieldMap(\Sigma, C, a) = \set{\bind{\seq{x}}{\seq{e}}} \\
  \Sigma; \Gamma \vdash \seq{e} \checks \seq{\tau}
}
{
  \Sigma; \Gamma \vdash \exConstr{q}{a} \synth C
}
\and
\inferrule*[lab={\ruleName{binop}}]
{
  \Sigma; \Gamma \vdash e \synth \sigma \\
  \Sigma; \Gamma \vdash e' \synth \sigma' \\
  \resolveOp(\Sigma, \odot, \sigma, \sigma') = \tau^\dagger
}
{
  \Sigma; \Gamma \vdash \exBinOp{e}{\odot}{e'} \synth \tau^\dagger
}
\and
\inferrule*[lab={\ruleName{unop}}]
{
  \Sigma; \Gamma \vdash e \synth \sigma \\
  \exUnOp{\ominus}{e} \neq \exUnOp{\pyinline{-}}{n} \quad (\forall n) \\
  \resolveOp(\Sigma, \ominus, \sigma) = \tau^\dagger
}
{
  \Sigma; \Gamma \vdash \exUnOp{\ominus}{e} \synth \tau^\dagger
}
\and
\inferrule*[lab={\ruleName{unop-literal}}]
{
  \strut
}
{
  \Sigma; \Gamma \vdash \exUnOp{\pyinline{-}}{n} \synth \tyLiteral{-n}
}
\and
\inferrule*[lab={\ruleName{syn-cond}}]
{
  \Sigma; \Gamma \vdash e \checks \tyBool \\
  I = \set{i \mid \Sigma; \Gamma \vdash e_i \synth \sigma_i} \neq \varnothing \\
  \Sigma; \Gamma \vdash e_j \checks \bigJoinTy{\Sigma}{(\sigma_i \mid i \in I)} \quad (\forall j \notin I)
}
{
  \Sigma; \Gamma \vdash \exCondExpr{e_1}{e}{e_2} \synth \bigJoinTy{\Sigma}{(\sigma_i \mid i \in I)}
}
\and
\inferrule*[lab={\ruleName{and}}]
{
  \Sigma; \Gamma \vdash e \checks \tyBool \\
  \Sigma; \Gamma \vdash e' \checks \tyBool
}
{
  \Sigma; \Gamma \vdash \exAnd{e}{e'} \synth \tyBool
}
\and
\inferrule*[lab={\ruleName{or}}]
{
  \Sigma; \Gamma \vdash e \checks \tyBool \\
  \Sigma; \Gamma \vdash e' \checks \tyBool
}
{
  \Sigma; \Gamma \vdash \exOr{e}{e'} \synth \tyBool
}
\end{mathpar}
\caption{Expression synthesis: names, literals, attributes, calls and operators}
\label{fig:typing-synth}
\end{figure}

\begin{figure}
\flushleft \shadebox{$\Sigma; \Gamma \vdash e \synth \tau$ (continued)}
\begin{mathpar}
\small
\inferrule*[lab={\ruleName{subscript-list}}]
{
  \Sigma; \Gamma \vdash e \synth \tyList{\tau} \\
  \Sigma; \Gamma \vdash e' \checks \tyInt
}
{
  \Sigma; \Gamma \vdash \exSubscript{e}{e'} \synth \tau
}
\and
\inferrule*[lab={\ruleName{subscript-str}}]
{
  \Sigma; \Gamma \vdash e \synth \tyStr \\
  \Sigma; \Gamma \vdash e' \checks \tyInt
}
{
  \Sigma; \Gamma \vdash \exSubscript{e}{e'} \synth \tyStr
}
\and
\inferrule*[lab={\ruleName{subscript-dict}}]
{
  \Sigma; \Gamma \vdash e \synth \tyDict{\tau} \\
  \Sigma; \Gamma \vdash e' \checks \tyStr
}
{
  \Sigma; \Gamma \vdash \exSubscript{e}{e'} \synth \tau
}
\and
\inferrule*[lab={\ruleName{subscript-tuple}}]
{
  \Sigma; \Gamma \vdash e \synth \tyTuple{\tau_1, \ldots, \tau_m} \\
  \Sigma; \Gamma \vdash e' \synth \tyLiteral{n} \\
  0 \leq n < m
}
{
  \Sigma; \Gamma \vdash \exSubscript{e}{e'} \synth \tau_{n+1}
}
\and
\inferrule*[lab={\ruleName{subscript-tuple-neg}}]
{
  \Sigma; \Gamma \vdash e \synth \tyTuple{\tau_1, \ldots, \tau_m} \\
  \Sigma; \Gamma \vdash e' \synth \tyLiteral{-n} \\
  \\
  0 < n \leq m
}
{
  \Sigma; \Gamma \vdash \exSubscript{e}{e'} \synth \tau_{(m-n)+1}
}
\and
\inferrule*[lab={\ruleName{subscript-tuple-int}}]
{
  \Sigma; \Gamma \vdash e \synth \tyTuple{\seq{\tau}} \\
  \Sigma; \Gamma \vdash e' \synth \tyInt
}
{
  \Sigma; \Gamma \vdash \exSubscript{e}{e'} \synth \bigJoinTy{\Sigma}{\seq{\tau}}
}
\and
\inferrule*[lab={\ruleName{subscript-union}}]
{
  \Sigma; \Gamma \vdash e \synth \tyUnion{\sigma}{\sigma'} \\
  \Sigma; \Gamma \override \set{\bind{y}{\sigma}} \vdash \exSubscript{y}{e'} \synth \tau \\
  \Sigma; \Gamma \override \set{\bind{y}{\sigma'}} \vdash \exSubscript{y}{e'} \synth \tau' \\
  y \notin \fv(e')
}
{
  \Sigma; \Gamma \vdash \exSubscript{e}{e'} \synth \joinTy{\Sigma}{\tau}{\tau'}
}
\and
\inferrule*[lab={\ruleName{attr-union}}]
{
  \Sigma; \Gamma \vdash e \synth \tyUnion{\sigma}{\sigma'} \\
  \Sigma; \Gamma \override \set{\bind{y}{\sigma}} \vdash \exAttr{y}{x} \synth \tau \\
  \Sigma; \Gamma \override \set{\bind{y}{\sigma'}} \vdash \exAttr{y}{x} \synth \tau'
}
{
  \Sigma; \Gamma \vdash \exAttr{e}{x} \synth \joinTy{\Sigma}{\tau}{\tau'}
}
\and
\inferrule*[lab={\ruleName{call-union}}]
{
  \Sigma; \Gamma \vdash e \synth \tyUnion{\sigma}{\sigma'} \\
  \Sigma; \Gamma \override \set{\bind{y}{\sigma}} \vdash \exCall{y}{\seq{e}} \synth \tau \\
  \Sigma; \Gamma \override \set{\bind{y}{\sigma'}} \vdash \exCall{y}{\seq{e}} \synth \tau' \\
  y \notin \fv(\seq{e})
}
{
  \Sigma; \Gamma \vdash \exCall{e}{\seq{e}} \synth \joinTy{\Sigma}{\tau}{\tau'}
}
\end{mathpar}
\caption{Expression synthesis: subscripts, and elimination at a union}
\label{fig:typing-synth-elim}
\end{figure}

\begin{figure}
\flushleft \shadebox{$\Sigma; \Gamma \vdash e \synth \tau$ (continued)}
\begin{mathpar}
\small
\inferrule*[lab={\ruleName{syn-list}}]
{
  I = \set{i \mid \Sigma; \Gamma \vdash e_i \synth \sigma_i} \neq \varnothing \\
  \tau = \bigJoinTy{\Sigma}{(\baseType(\sigma_i) \mid i \in I)} \\
  \Sigma; \Gamma \vdash e_j \checks \tau \quad (\forall j \notin I)
}
{
  \Sigma; \Gamma \vdash \exList{e_1, \ldots, e_{n+1}} \synth \tyList{\tau}
}
\and
\inferrule*[lab={\ruleName{syn-tuple}}]
{
  \Sigma; \Gamma \vdash e_i \synth \tau_i \quad (\forall i)
}
{
  \Sigma; \Gamma \vdash \exTuple{\seq{e}} \synth \tyTuple{\seq{\tau}}
}
\and
\inferrule*[lab={\ruleName{syn-dict}}]
{
  \Sigma; \Gamma \vdash e_i \checks \tyStr \quad (\forall i) \\
  I = \set{i \mid \Sigma; \Gamma \vdash e_i' \synth \sigma_i} \neq \varnothing \\
  \tau = \bigJoinTy{\Sigma}{(\baseType(\sigma_i) \mid i \in I)} \\
  \Sigma; \Gamma \vdash e_j' \checks \tau \quad (\forall j \notin I)
}
{
  \Sigma; \Gamma \vdash \exDict{\bind{e_1}{e_1'}, \ldots, \bind{e_{n+1}}{e_{n+1}'}} \synth \tyDict{\tau}
}
\and
\inferrule*[lab={\ruleName{syn-list-comp}}]
{
  \Sigma; \Gamma \vdash \seq{g} : \Delta \\
  \Sigma; \Gamma \override \Delta \vdash e \synth \tau \\
  \captures(e) \cap \dom(\Delta) = \varnothing
}
{
  \Sigma; \Gamma \vdash \exListComp{e}{\seq{g}} \synth \tyList{\baseType(\tau)}
}
\and
\inferrule*[lab={\ruleName{syn-dict-comp}}]
{
  \Sigma; \Gamma \vdash \seq{g} : \Delta \\
  \Sigma; \Gamma \override \Delta \vdash e \checks \tyStr \\
  \Sigma; \Gamma \override \Delta \vdash e' \synth \tau \\
  \captures((e, e')) \cap \dom(\Delta) = \varnothing
}
{
  \Sigma; \Gamma \vdash \exDictComp{e}{e'}{\seq{g}} \synth \tyDict{\baseType(\tau)}
}
\end{mathpar}
\caption{Expression synthesis: tuples, lists, dictionaries and comprehensions}
\label{fig:typing-synth-ctd}
\end{figure}


The fields of a class $C$ form a sequence $\fields(\Sigma, C)$ of bindings $\bind{x}{\tau}$, one per
field: the fields of its base class first, recursively, then the fields $C$ declares itself, each at its
declared type. The order matters, because a constructor call pairs positional arguments with fields by
position. A constructor call $\exConstr{q}{a}$ pairs each argument with a field using
$\fieldMap$, checks it against that field's type, and synthesises the class $C$ that $q$ resolves to
(\ruleName{constr}). In an attribute reference $\exAttr{e}{x}$ the object $e$ must
synthesise a class $C$, and $\exAttr{e}{x}$ then has the type bound to $x$ in $\fields(\Sigma, C)$
(\ruleName{attr-object}). The signature of a predefined module binds
$\pymath{\_\_name\_\_}$ at $\tyStr$ and each member at the type listed for it
(\figref{predefined-types}, Definition~\ref{def:predefined}). A module attribute reference
$\exAttr{e}{x}$ resolves as a qualified name to a definitely assigned member
(\ruleName{attr-module})\paperonly{; see \secref{no-load-dependent}}. A variable bound to a module or
class is not an expression on its own, and occurs only as the head of a qualified
name\paperonly{ (\secref{no-first-class})}.
In a subscript $\exSubscript{e}{e'}$ there is a rule for each container $e$ can synthesise: a list
(\ruleName{subscript-list}) or string (\ruleName{subscript-str}) with $e'$ of type \pyinline{int}, and a
dictionary with $e'$ of type \pyinline{str} (\ruleName{subscript-dict}). Tuples have three rules: an index
that synthesises a literal type gives the component type at that position, counting from the end where the index
is negative (\ruleName{subscript-tuple}, \ruleName{subscript-tuple-neg}), and any other index of type
\pyinline{int} gives the join of the component types (\ruleName{subscript-tuple-int}). A negated numeric
literal synthesises the negated literal type (\ruleName{unop-literal}), and \ruleName{unop} excludes that
form, to keep the judgement deterministic. A subscript, attribute reference or call whose subject
synthesises a union is typed at each member, with a variable in place of the subject, and synthesises the
join of the results (\ruleName{subscript-union}, \ruleName{attr-union}, \ruleName{call-union}); generators need no
such rule, because $\elemTy$ takes a union to the join of its members' element types.
\citet{dunfield2004} state union elimination once, for an arbitrary evaluation context; here each
elimination form has its own rule.

Lambdas check, as do lists, tuples, dictionaries and comprehensions. An expression that synthesises $\sigma$ checks against any supertype
of $\sigma$ (\ruleName{check-synth}), which applies only where no other rule concludes $\Sigma; \Gamma \vdash e \checks
\tau$, so both judgements are
syntax-directed and an expression has at most one derivation in each.

\begin{figure}
\begin{subfigure}{\textwidth}
\flushleft \shadebox{$\Sigma; \Gamma \vdash e \checks \tau$}
\begin{mathpar}
\small
\inferrule*[lab={\ruleName{check-synth}}]
{
  \Sigma; \Gamma \vdash e \synth \sigma \\
  \Sigma \vdash \sigma \subty \tau \\
  \text{no other rule concludes } \Sigma; \Gamma \vdash e \checks \tau
}
{
  \Sigma; \Gamma \vdash e \checks \tau
}
\and
\inferrule*[lab={\ruleName{lambda}}]
{
  \Sigma; \Gamma \override \set{\bind{x_1}{\tau_1}, \ldots, \bind{x_n}{\tau_n}} \vdash e \checks \tau
}
{
  \Sigma; \Gamma \vdash \exLambda{x_1, \ldots, x_n}{e} \checks \tyCallable{\tau_1, \ldots, \tau_n}{\tau}
}
\and
\inferrule*[lab={\ruleName{call-lambda}}]
{
  \Sigma; \Gamma \vdash e_i \synth \tau_i \quad (\forall i) \\
  \Sigma; \Gamma \override \set{\bind{x_1}{\tau_1}, \ldots, \bind{x_n}{\tau_n}} \vdash e \checks \tau
}
{
  \Sigma; \Gamma \vdash \exCall{(\exLambda{x_1, \ldots, x_n}{e})}{e_1, \ldots, e_n} \checks \tau
}
\and
\inferrule*[lab={\ruleName{list}}]
{
  \Sigma; \Gamma \vdash e_i \checks \tau \quad (\forall i)
}
{
  \Sigma; \Gamma \vdash \exList{e_1, \ldots, e_n} \checks \tyList{\tau}
}
\and
\inferrule*[lab={\ruleName{tuple}}]
{
  \Sigma; \Gamma \vdash e_i \checks \tau_i \quad (\forall i)
}
{
  \Sigma; \Gamma \vdash \exTuple{e_1, \ldots, e_n} \checks \tyTuple{\tau_1, \ldots, \tau_n}
}
\and
\inferrule*[lab={\ruleName{dict}}]
{
  \Sigma; \Gamma \vdash e_i \checks \tyStr \quad (\forall i) \\
  \Sigma; \Gamma \vdash e_i' \checks \tau \quad (\forall i)
}
{
  \Sigma; \Gamma \vdash \exDict{\bind{e_1}{e_1'}, \ldots, \bind{e_n}{e_n'}} \checks \tyDict{\tau}
}
\and
\inferrule*[lab={\ruleName{cond}}]
{
  \Sigma; \Gamma \vdash e \checks \tyBool \\
  \Sigma; \Gamma \vdash e_1 \checks \tau \\
  \Sigma; \Gamma \vdash e_2 \checks \tau
}
{
  \Sigma; \Gamma \vdash \exCondExpr{e_1}{e}{e_2} \checks \tau
}
\and
\inferrule*[lab={\ruleName{list-comp}}]
{
  \Sigma; \Gamma \vdash \seq{g} : \Delta \\
  \Sigma; \Gamma \override \Delta \vdash e \checks \tau \\
  \capturesE(e) \cap \dom(\Delta) = \varnothing
}
{
  \Sigma; \Gamma \vdash \exListComp{e}{g_1\;\ldots\;g_n} \checks \tyList{\tau}
}
\and
\inferrule*[lab={\ruleName{dict-comp}}]
{
  \Sigma; \Gamma \vdash \seq{g} : \Delta \\
  \Sigma; \Gamma \override \Delta \vdash e \checks \tyStr \\
  \Sigma; \Gamma \override \Delta \vdash e' \checks \tau \\
  \capturesE((e, e')) \cap \dom(\Delta) = \varnothing
}
{
  \Sigma; \Gamma \vdash \exDictComp{e}{e'}{g_1\;\ldots\;g_n} \checks \tyDict{\tau}
}
\end{mathpar}
\subcaption{Expression $e$ checks against type $\tau$}
\end{subfigure}

\vspace{2mm}

\begin{subfigure}{\textwidth}
\flushleft \shadebox{$\Sigma; \Gamma \vdash \seq{g} : \Delta$}
\begin{mathpar}
\small
\inferrule*[lab={\ruleName{quals-nil}}]
{
  \strut
}
{
  \Sigma; \Gamma \vdash \varepsilon : \varnothing
}
\and
\inferrule*[lab={\ruleName{qual-guard}}]
{
  \Sigma; \Gamma \vdash e \checks \tyBool \\
  \Sigma; \Gamma \vdash \seq{g} : \Delta
}
{
  \Sigma; \Gamma \vdash \qualGuard{e} \cons \seq{g} : \Delta
}
\and
\inferrule*[lab={\ruleName{qual-generator}}]
{
  \Sigma; \Gamma \vdash e \synth \tau \\
  \elemTy(\Sigma, \tau) = \sigma \\
  \Sigma; \Gamma \override \set{\bind{x}{\sigma}} \vdash \seq{g} : \Delta \\
  x \notin \capturesS(\seq{g})
}
{
  \Sigma; \Gamma \vdash \qualGenerator{x}{e} \cons \seq{g} : \set{\bind{x}{\sigma}} \override \Delta
}
\end{mathpar}
\subcaption{Comprehension qualifiers $\seq{g}$ bind $\Delta$}
\end{subfigure}
\caption{Expression checking and comprehension qualifiers}
\label{fig:typing-check}
\end{figure}


The forms with synthesis rules of their own are conditional expressions when a branch
synthesises, at the join of the two types when both branches synthesise and otherwise at the
synthesising branch's type, the other branch checking against it (\ruleName{syn-cond});
tuples when every component synthesises (\ruleName{syn-tuple}); comprehensions when the body
synthesises (\ruleName{syn-list-comp}, \ruleName{syn-dict-comp}); and non-empty lists and
dictionaries when at least one element synthesises, the remaining elements checking against the join
of the base types of the synthesising elements (\ruleName{syn-list}, \ruleName{syn-dict}). A list,
dictionary or comprehension takes the base type of the type its elements or body synthesise, since
\pyinline{list} and \pyinline{dict} are invariant and a literal element type would keep the container from
matching any other element type. An empty list or dictionary and a lambda have no
synthesis rule at all, so checking is the only way to type them. A lambda applied to arguments is the
exception, since the arguments give the parameter types, and the call then synthesises the type of the body
or checks the body against an expected type (\ruleName{syn-call-lambda}, \ruleName{call-lambda}). A non-empty list has both, and the
checking rule is still needed where the expected element type is wider than the one the elements
synthesise, since \pyinline{list} is invariant.

Typing a comprehension types its qualifiers, which bind the variables that can be used in the body
(\figref{typing-check}). A generator's variable is bound at $\elemTy$ of the type synthesised by the
iterated expression (\secref{types-auxiliary}). A
tuple is iterated at the join of the base types of its components.
A lambda within a comprehension may not capture a generator variable, a side condition of
\ruleName{list-comp}, \ruleName{dict-comp} and \ruleName{qual-generator}. Python closures are
late-binding and the comprehension rebinds the variable on each iteration, so the lambda would see the
variable's final value; in \PurePy a lambda is only able to see the value of the generator variable at
each iteration.


\speconly{\subsubsection{Operators}\label{sec:well-formedness-operators}}%
\paperonly{\paragraph{Operators.}}

Each operator is overloaded as specified in \figref{operator-overloads}. A \emph{resolved overload}
$((\tau_1, \tau_2), \tau)$ pairs operand \emph{bounds}
with a result type; the judgements of the figure relate operand
types to the resolved overloads of an operator.

\begin{figure}
\small
\setlength{\jot}{2pt}
\flushleft \shadebox{$\Sigma \vdash \odot(\sigma_1, \sigma_2) : ((\tau_1, \tau_2), \tau)$}
\begin{defalign}
& \Sigma \vdash \odot(\sigma_1, \sigma_2) : ((\sigma_1, \sigma_2), \tyBool) && \quad \text{if } \odot \in \set{\pyinline{==}, \pyinline{!=}} \text{ and } \sigma_1, \sigma_2 \text{ comparable equality types} \\
& \Sigma \vdash \odot(\sigma_1, \tyList{\tau}) : ((\sigma_1, \tyList{\tau}), \tyBool) && \quad \text{if } \odot \in \set{\pyinline{in}, \pyinline{not in}} \text{ and } \sigma_1, \tau \text{ comparable} \\
& \Sigma \vdash \odot(\sigma_1, \tyTuple{\seq{\tau}}) : ((\sigma_1, \tyTuple{\seq{\tau}}), \tyBool) && \quad \text{if } \odot \in \set{\pyinline{in}, \pyinline{not in}} \text{ and } \sigma_1, \bigJoinTy{\Sigma}{\seq{\tau}} \text{ comparable} \\
& \Sigma \vdash \odot(\sigma_1, \sigma_2) : ((\tyStr, \tyStr), \tyBool) && \quad \text{if } \odot \in \set{\pyinline{in}, \pyinline{not in}} \text{ and } \Sigma \vdash \sigma_1, \sigma_2 \subty \tyStr \\
& \Sigma \vdash \odot(\sigma_1, \tyDict{\tau}) : ((\tyStr, \tyDict{\tau}), \tyBool) && \quad \text{if } \odot \in \set{\pyinline{in}, \pyinline{not in}} \text{ and } \Sigma \vdash \sigma_1 \subty \tyStr \\
& \Sigma \vdash \odot(\sigma_1, \sigma_2) : ((\tyFloat, \tyFloat), \tyBool) && \quad \text{if } \odot \in \set{\pyinline{<}, \pyinline{<=}, \pyinline{>}, \pyinline{>=}} \text{ and } \Sigma \vdash \sigma_1, \sigma_2 \subty \tyFloat \\
& \Sigma \vdash \odot(\sigma_1, \sigma_2) : ((\tyStr, \tyStr), \tyBool) && \quad \text{if } \odot \in \set{\pyinline{<}, \pyinline{<=}, \pyinline{>}, \pyinline{>=}} \text{ and } \Sigma \vdash \sigma_1, \sigma_2 \subty \tyStr \\
& \Sigma \vdash \odot(\sigma_1, \sigma_2) : ((\tyInt, \tyInt), \tyInt) && \quad \text{if } \odot \in \set{\pyinline{+}, \pyinline{-}, \pyinline{*}, \pyinline{//}, \pyinline{\%}} \text{ and } \Sigma \vdash \sigma_1, \sigma_2 \subty \tyInt \\
& \Sigma \vdash \odot(\sigma_1, \sigma_2) : ((\tyFloat, \tyFloat), \tyFloat) && \quad \text{if } \odot \in \set{\pyinline{+}, \pyinline{-}, \pyinline{*}, \pyinline{//}, \pyinline{\%}, \pyinline{/}, \pyinline{**}} \text{ and } \Sigma \vdash \sigma_1, \sigma_2 \subty \tyFloat \\
& \Sigma \vdash \pyinline{+}(\sigma_1, \sigma_2) : ((\tyStr, \tyStr), \tyStr) && \quad \text{if } \Sigma \vdash \sigma_1, \sigma_2 \subty \tyStr \\
& \rlap{$\Sigma \vdash \pyinline{+}(\tyList{\sigma}, \tyList{\tau}) : ((\tyList{\sigma}, \tyList{\tau}), \tyList{\joinTy{\Sigma}{\sigma}{\tau}})$} \\
& \rlap{$\Sigma \vdash \pyinline{+}(\tyTuple{\seq{\sigma}}, \tyTuple{\seq{\tau}}) : ((\tyTuple{\seq{\sigma}}, \tyTuple{\seq{\tau}}), \tyTuple{\seq{\sigma} \mathbin{\concat} \seq{\tau}})$} \\
& \Sigma \vdash \pyinline{*}(\sigma_1, \sigma_2) : ((\tyStr, \tyInt), \tyStr) && \quad \text{if } \Sigma \vdash \sigma_1 \subty \tyStr \text{ and } \Sigma \vdash \sigma_2 \subty \tyInt \\
& \Sigma \vdash \pyinline{*}(\sigma_1, \sigma_2) : ((\tyInt, \tyStr), \tyStr) && \quad \text{if } \Sigma \vdash \sigma_1 \subty \tyInt \text{ and } \Sigma \vdash \sigma_2 \subty \tyStr \\
& \Sigma \vdash \pyinline{*}(\tyList{\tau}, \sigma_2) : ((\tyList{\tau}, \tyInt), \tyList{\tau}) && \quad \text{if } \Sigma \vdash \sigma_2 \subty \tyInt \\
& \Sigma \vdash \pyinline{*}(\sigma_1, \tyList{\tau}) : ((\tyInt, \tyList{\tau}), \tyList{\tau}) && \quad \text{if } \Sigma \vdash \sigma_1 \subty \tyInt \\
& \Sigma \vdash \pyinline{**}(\sigma_1, \tyLiteral{n}) : ((\tyInt, \tyLiteral{n}), \tyInt) && \quad \text{if } \Sigma \vdash \sigma_1 \subty \tyInt \text{ and } n \text{ an integer} \\
& \Sigma \vdash \pyinline{**}(\sigma_1, \tyLiteral{-n}) : ((\tyInt, \tyLiteral{-n}), \tyFloat) && \quad \text{if } \Sigma \vdash \sigma_1 \subty \tyInt \text{ and } n \text{ a positive integer} \\
& \Sigma \vdash \odot(\sigma_1, \sigma_2) : ((\tau_1, \tau_2), \tau) && \quad \text{if } \Sigma \vdash {\odot}(\baseType(\sigma_1), \baseType(\sigma_2)) : ((\tau_1, \tau_2), \tau)
\end{defalign}

\vspace{-2mm}

\flushleft \shadebox{$\Sigma \vdash \ominus(\sigma) : (\tau, \tau')$}
\begin{defalign}
& \pyinline{not}(\sigma) : (\tyBool, \tyBool) && \quad \text{if } \Sigma \vdash \sigma \subty \tyBool \\
& \Sigma \vdash \ominus(\sigma) : (\tyInt, \tyInt) && \quad \text{if } \ominus \in \set{\pyinline{+}, \pyinline{-}} \text{ and } \Sigma \vdash \sigma \subty \tyInt \\
& \Sigma \vdash \ominus(\sigma) : (\tyFloat, \tyFloat) && \quad \text{if } \ominus \in \set{\pyinline{+}, \pyinline{-}} \text{ and } \Sigma \vdash \sigma \subty \tyFloat \\
& \Sigma \vdash \ominus(\sigma) : (\tau, \tau') && \quad \text{if } \Sigma \vdash {\ominus}(\baseType(\sigma)) : (\tau, \tau')
\end{defalign}
\caption{Operator overloads}
\label{fig:operator-overloads}
\end{figure}


\begin{definition}[Overloads at operand types]
\label{def:overloads}
The set $\overloadsOf(\Sigma, \odot, \sigma_1, \sigma_2)$ consists of the tuples $((\tau_1, \tau_2), \tau)$ such
that $\Sigma \vdash \odot(\sigma_1, \sigma_2) : ((\tau_1, \tau_2), \tau)$, ordered by componentwise subtyping in $\Sigma$ on their
bounds $(\tau_1, \tau_2)$. The unary set $\overloadsOf(\Sigma, \ominus, \sigma)$ is defined in the same way from
conclusions $\Sigma \vdash \ominus(\sigma) : (\tau, \tau')$. In both cases $\min$ denotes the least element
under this order, where one exists.
\end{definition}

\noindent The \ruleName{binop} and \ruleName{unop} rules (\figref{typing-synth}) take the result type
of $\min(\overloadsOf(\Sigma, \odot, \sigma_1, \sigma_2))$, so \pyinline{1 + 2} is an
\pyinline{int} and \pyinline{1 + 2.0} a \pyinline{float}.

An \emph{equality type} is one for which structural equality is
defined:
\begin{defalign}
& \Sigma \vdash \nu \eqType \\
& \Sigma \vdash \tyLiteral{\ell} \eqType \\
& \Sigma \vdash \tyList{\tau} \eqType && \qquad\quad \text{if } \Sigma \vdash \tau \eqType \\
& \Sigma \vdash \tyDict{\tau} \eqType && \qquad\quad \text{if } \Sigma \vdash \tau \eqType \\
& \Sigma \vdash \tyTuple{\seq{\tau}} \eqType && \qquad\quad \text{if } \Sigma \vdash \tau_i \eqType \quad (\forall i) \\
& \Sigma \vdash \tyUnion{\sigma}{\tau} \eqType && \qquad\quad \text{if } \Sigma \vdash \sigma \eqType \text{ and } \Sigma \vdash \tau \eqType \\
& \Sigma \vdash C \eqType && \qquad\quad \text{if } \Sigma \vdash \tau \eqType \text{ for each } \bind{x}{\tau} \in \fields(\Sigma, C)
\end{defalign}

\noindent A callable type is not an equality type, and neither is any type with a callable component, so
comparing two functions is rejected statically. The final rule of each table closes the
judgement under base-typing, letting \pyinline{1 == 0} use the equality overload at \pyinline{int}, since
two distinct literal types are not comparable; \pyinline{2 ** 3} is still an \pyinline{int}, since the
direct instance's bounds $(\tyInt, \tyLiteral{3})$ lie below the base instance's
$(\tyFloat, \tyFloat)$. Resolution is more precise than mypy for this operator, whose overloads give
\pyinline{Any} as soon as the exponent is not a literal, where resolution gives \pyinline{float}.

\begin{lemma}[Least overload]
\label{lem:least-overload}
If neither $\sigma_1$ nor $\sigma_2$ is a subtype of \pyinline{Never} and
$\overloadsOf(\Sigma, \odot, \sigma_1, \sigma_2)$ is not empty, then the set has a least element; likewise
$\overloadsOf(\Sigma, \ominus, \sigma)$.
\end{lemma}

\subsection{Statements}
\label{sec:well-formedness-statements}

Every well-typed statement either \emph{definitely returns}, or \emph{definitely assigns} a context
$\Delta$, which may be empty. The base cases are that \pyinline{return e} definitely returns, and
\pyinline{x = e} definitely assigns $\set{\bind{\pyinline{x}}{\sigma}}$ for the type $\sigma$ of
\pyinline{e}. \emph{Static outcomes} $R$ classify the two completions: $R$ is either $\Returns$ (definitely
returns) or $\Assigns{\Delta}$ (definitely assigns $\Delta$).

\begin{figure}
\flushleft \shadebox{$\Sigma; \Gamma; \zeta \vdash s : R$}
\begin{mathpar}
\small
\inferrule*[lab={\ruleName{pass}}]
{
  \strut
}
{
  \Sigma; \Gamma; \zeta \vdash \exPass : \Assigns{\varnothing}
}
\and
\inferrule*[lab={\ruleName{declare}}]
{
  \Gamma \vdash \psi \synth \tau
}
{
  \Sigma; \Gamma; \zeta \vdash x\exAnnot{\psi} : \Assigns{\set{\bind{x}{\DU{\tau}}}}
}
\and
\inferrule*[lab={\ruleName{assign}}]
{
  \Gamma(x) = \DU{\tau} \\
  \Sigma; \Gamma \vdash e \checks \tau
}
{
  \Sigma; \Gamma; \zeta \vdash \exAssign{x}{e} : \Assigns{\set{\bind{x}{\tau}}}
}
\and
\inferrule*[lab={\ruleName{declare-assign}}]
{
  \Gamma \vdash \psi \synth \tau \\
  \Sigma; \Gamma \vdash e \checks \tau
}
{
  \Sigma; \Gamma; \zeta \vdash \exAssign{x\exAnnot{\psi}}{e} : \Assigns{\set{\bind{x}{\tau}}}
}
\and
\inferrule*[lab={\ruleName{expr-stmt}}]
{
  \Sigma; \Gamma \vdash e \synth \sigma
}
{
  \Sigma; \Gamma; \zeta \vdash e : \Assigns{\varnothing}
}
\and
\inferrule*[lab={\ruleName{assert}}]
{
  \Sigma; \Gamma \vdash e \checks \tyBool
}
{
  \Sigma; \Gamma; \zeta \vdash \exAssert{e} : \Assigns{\varnothing}
}
\and
\inferrule*[lab={\ruleName{assert-msg}}]
{
  \Sigma; \Gamma \vdash e \checks \tyBool \\
  \Sigma; \Gamma \vdash e' \checks \tyStr
}
{
  \Sigma; \Gamma; \zeta \vdash \exAssertMsg{e}{e'} : \Assigns{\varnothing}
}
\and
\inferrule*[lab={\ruleName{return-none}}]
{
  \Sigma \vdash \tyNone \subty \tau
}
{
  \Sigma; \Gamma; \tau \vdash \exReturnNone : \Returns
}
\and
\inferrule*[lab={\ruleName{return}}]
{
  \Sigma; \Gamma \vdash e \checks \tau
}
{
  \Sigma; \Gamma; \tau \vdash \exReturn{e} : \Returns
}
\and
\inferrule*[lab={\ruleName{if}}]
{
  \Sigma; \Gamma \vdash e_i \checks \tyBool \quad (\forall i) \\
  \Sigma; \Gamma; \zeta \vdash s_i : R_i \quad (\forall i)
}
{
  \Sigma; \Gamma; \zeta \vdash \exIf{e_1}{s_1}{\exElif{e_2}{s_2}\;\ldots\;\exElif{e_{n+1}}{s_{n+1}}} : R_1 \merge \ldots \merge R_{n+1} \: \merge \Assigns{\varnothing}
}
\and
\inferrule*[lab={\ruleName{if-else}}]
{
  \Sigma; \Gamma \vdash e_i \checks \tyBool \quad (\forall i) \\
  \Sigma; \Gamma; \zeta \vdash s_i : R_i \quad (\forall i)
}
{
  \Sigma; \Gamma; \zeta \vdash \exIfElse{e_1}{s_1}{\exElif{e_2}{s_2}\;\ldots\;\exElif{e_{n+1}}{s_{n+1}}}{s_{n+2}} : R_1 \merge \ldots \merge R_{n+2}
}
\and
\inferrule*[lab={\ruleName{seq}}]
{
  \Sigma; \Gamma; \zeta \vdash s : \Assigns{\Delta} \\
  \Sigma; \Gamma \override \Delta ; \zeta \vdash s' : R
}
{
  \Sigma; \Gamma; \zeta \vdash s\;s' : \Assigns{\Delta} \override R
}
\and
\inferrule*[lab={\ruleName{def}}]
{
  \Gamma \vdash d_i \synth \tau_i \quad (\forall i) \\
  \Delta' = \set{\bind{\declaresS(\seq{d})}{\seq{\tau}}} \\
  \Sigma; \Gamma \override \Delta' \vdash d_i \checks \tau_i \quad (\forall i)
}
{
  \Sigma; \Gamma; \zeta \vdash \seq{d} : \Assigns{\Delta'}
}
\and
\inferrule*[lab={\ruleName{match}}]
{
  \Sigma; \Gamma \vdash e \synth \sigma \\
  \Sigma; \Gamma; \zeta; \sigma \vdash \shapes(\Sigma, \sigma, \varnothing), \seq{b} : S, \varnothing
}
{
  \Sigma; \Gamma; \zeta \vdash \exMatch{e}{\seq{b}} : S
}
\and
\inferrule*[lab={\ruleName{match-partial}}]
{
  \Sigma; \Gamma \vdash e \synth \sigma \\
  \Sigma; \Gamma; \zeta; \sigma \vdash \shapes(\Sigma, \sigma, \varnothing), \seq{b} : S, L \\
  L \neq \varnothing
}
{
  \Sigma; \Gamma; \zeta \vdash \exMatch{e}{\seq{b}} : S \merge \Assigns{\varnothing}
}
\end{mathpar}
\caption{Given optional return type $\zeta$, statement $s$ has static outcome $R$}
\label{fig:typing-stmt}
\end{figure}


\begin{figure}
\flushleft \shadebox{$\Sigma; \Gamma; \zeta; \sigma \vdash \seq{h} : S, L$}
\begin{mathpar}
\small
\inferrule*[lab={\ruleName{cases}}]
{
  L_0 = \shapesOf(\Sigma, \sigma, \varnothing) \\
  \Sigma; \Gamma \vdash p_i \agrees \sigma \quad (\forall i) \\
  \Sigma; \Gamma \vdash L_{i-1}, p_i \matches K_i, L_i \binding \Delta_i \quad (\forall i) \\
  \Sigma; \Gamma \override \Delta_i; \zeta \vdash s_i : R_i \quad (\forall i) \\
  S_i = \Assigns{\Delta_i} \override R_i \quad (\forall i)
}
{
  \Sigma; \Gamma; \zeta; \sigma \vdash \exCase{p_1}{s_1}\;\ldots\;\exCase{p_{n+1}}{s_{n+1}} : S_1 \mergeS \ldots \mergeS S_{n+1}, L_{n+1}
}

\end{mathpar}
\caption{Given optional return type $\zeta$ and scrutinee type $\sigma$, cases $\seq{h}$ have static outcome $S$ and leave residual $L$}
\label{fig:typing-cases}
\end{figure}


\begin{figure}
\flushleft \shadebox{$\Sigma; q; \Gamma \vdash t : R \binding \Sigma'$}
\begin{mathpar}
\small
\inferrule*[lab={\ruleName{top-stmt}}]
{
  \Sigma; \Gamma; \varepsilon \vdash s : R
}
{
  \Sigma; q; \Gamma \vdash s : R \binding \Sigma
}
\and
\inferrule*[lab={\ruleName{top-seq}}]
{
  \Sigma; q; \Gamma \vdash t : \Assigns{\Delta} \binding \Sigma' \\
  \Sigma'; q; \Gamma \override \Delta \vdash t' : R \binding \Sigma^\dagger \\
  \capturesS(t) \cap \assignsS(t') = \varnothing \\
  \nexists c, C.\ \bind{c}{\Type{C}} \in \Delta \text{ and } c \in \assignsS(t')
}
{
  \Sigma; q; \Gamma \vdash t\;t' : \Assigns{\Delta} \override R \binding \Sigma^\dagger
}
\and
\inferrule*[lab={\ruleName{class}}]
{
  \bind{\pyinline{dataclass}}{\predefName} \in \Gamma \\
  \ownFields(\phi) = \bind{\seq{x}}{\seq{\psi}}, \; \seq{x} \text{ distinct} \\
  \Gamma \vdash \psi_i \synth \sigma_i \quad (\forall i)
}
{
  \Sigma; q; \Gamma \vdash \exDataclass{c}{\phi} : \Assigns{\set{\bind{c}{\Type{\fq{q.c}}}}} \binding \Sigma, \bind{\fq{q.c}}{\clsRec{\bind{\seq{x}}{\seq{\sigma}}}{\bot}}
}
\and
\inferrule*[lab={\ruleName{class-extend}}]
{
  \bind{c'}{\Type{D}} \in \Gamma \\
  \bind{\pyinline{dataclass}}{\predefName} \in \Gamma \\
  \ownFields(\phi) = \bind{\seq{x}}{\seq{\psi}}, \; \seq{x} \text{ distinct} \\
  \Gamma \vdash \psi_i \synth \sigma_i \quad (\forall i) \\
  \seq{x} \cap \dom(\fields(\Sigma, D)) = \varnothing
}
{
  \Sigma; q; \Gamma \vdash \exDataclassExtend{c}{c'}{\phi} : \Assigns{\set{\bind{c}{\Type{\fq{q.c}}}}} \binding \Sigma, \bind{\fq{q.c}}{\clsRec{\bind{\seq{x}}{\seq{\sigma}}}{D}}
}
\end{mathpar}
\caption{Top-level statement $t$ has result type $R$, extending $\Sigma$ to $\Sigma'$}
\label{fig:typing-top}
\end{figure}


Context operations $\override$ and $\mergeS$ lift to a right semiring over static outcomes $R$, with $\Assigns{\varnothing}$ as left and right unit for $\override$ and $\Returns$ as unit for $\mergeS$ and zero for $\override$.

\begin{definition}[Override and merge on static outcomes]
\begin{defalign}
\Assigns{\Delta} \mergeS \Assigns{\Delta'} & = \Assigns{\Delta \mergeS \Delta'} \\
R \mergeS \Returns = \Returns \mergeS R & = R \\
\\
\Assigns{\Delta} \override \Assigns{\Delta'} & = \Assigns{\Delta \override \Delta'} \\
R \override \Returns = \Returns \override R & = \Returns
\end{defalign}
\end{definition}

\noindent The $\Returns \override R$ case never arises because statement typing excludes unreachable statements\paperonly{ (\secref{unreachable})}.

\paperinput{paper/well-formedness/cases}%

The statement judgement $\Sigma; \Gamma; \zeta \vdash s : R$ says that under $\Gamma$ and
optional return type $\zeta$, statement $s$ has static outcome $R$ (\figref{typing-stmt}); every rule passes $\zeta$ down unchanged.
A sequence $s\;s'$ (\ruleName{seq}) types $s$, types $s'$ under $\Gamma$ overridden by the context assigned
by $s$, and requires that no variable captured by $s$ is assigned in $s'$. A conditional combines its
branches with $\mergeS$, merging in $\Assigns{\varnothing}$ when there is no \pyinline{else} branch.
Top-level statements have a judgement of
their own, $\Sigma; q; \Gamma \vdash t : R \binding \Sigma'$ (\figref{typing-top}): a sequence composes in the same way as a statement sequence, a
class declaration is a top-level form, and a plain statement is checked with no return type
(\ruleName{top-stmt}), so a \pyinline{return} has no derivation at the top level. A top-level
statement may not assign a name that an earlier top-level statement bound to a class
(\ruleName{top-seq}), so a class binding is permanent in its module and no later declaration or
assignment can take a declared class's name. A
\pyinline{return} checks its expression against $\zeta$, which must be of the form $\tau$, and a bare
\pyinline{return} requires $\tau$ to admit \pyinline{None}. An annotated assignment checks the expression against the type of the annotation,
and an unannotated assignment takes the type synthesised by the expression (\ruleName{assign},
\ruleName{assign-unannot}). The conditions of \pyinline{if} and \pyinline{assert}
check against \pyinline{bool}, and an assertion message against \pyinline{str}. A class declaration binds
the class name to $\Type{C}$, requiring the class's own fields to be distinct and, where it extends
another class, the base name to be bound to a class and the additional fields not to repeat an inherited field
(\ruleName{class}, \ruleName{class-extend}). Class declarations also require the predefined name \pyinline{dataclass} (\figref{predefined-types}) to be in
scope; every \PurePy class must declare itself a dataclass. Each rule that takes an annotation resolves it to a type
(\figref{resolve-type}).

A function definition binds its parameters at their annotated types, plus the names of the
mutual region it belongs to, including its own name, each at the appropriate $\pyinline{Callable}$ type; variables assigned in the body start
at $\False$ (\ruleName{def}), emulating Python's scoping of function bodies, which
prefers a locally assigned name to one visible in an enclosing scope. The body is checked with the
declared return type, and must have static outcome $\Returns$ unless the function returns \pyinline{None},
so a function that declares a result definitely returns one.

A \pyinline{match} scrutinee synthesises a type $\sigma$, and its cases $\seq{b}$ are checked using the auxiliary
judgement $\Sigma; \Gamma; \zeta; \sigma \vdash L, \seq{b} : S, L'$ (\figref{typing-cases}), which takes
the cases in turn: each case's pattern is matched against the \emph{residual} $L$, a static
over-approximation of the values of $\sigma$ which remain unmatched by earlier cases, its body is
checked under the pattern's bindings, and the outcomes
are merged into $S$. Each pattern must be sequence-safe at $\sigma$ (\secref{types-auxiliary}). The match rules start the
residual from the shapes of $\sigma$ (\figref{shapes}); \ruleName{match} applies where the final residual is empty; \ruleName{match-partial} applies otherwise, merging in
$\Assigns{\varnothing}$ for the values that fall through. Shapes, matching a pattern against the residual,
and sequence-safety are the subject of \secref{well-formedness-patterns}.

\subsection{Pattern matching}
\label{sec:exhaustiveness}
\label{sec:well-formedness-patterns}

To check a \pyinline{match}, each case $\exCase{p}{s}$ is matched in turn against the residual, an
ordered set of \emph{shapes} (\figref{shapes}). A shape denotes a set of values of a single type, and the
residual denotes the values of the scrutinee's type which remain unmatched by earlier cases. Matching the case removes from the residual the values matched
by $p$ (\figref{pat}); a case that does not match any shape of the residual is unreachable and the match is rejected.
The match is exhaustive if the last case leaves the residual empty, and may
otherwise fall through (\ruleName{match}, \ruleName{match-partial}). Checking reachability and
exhaustiveness with one relation follows the usefulness algorithm for ML~\cite{maranget2007}. The residual plays the
role of the term descriptions of \citet{sestoft1996}, and the uncovered sets of \citet{karachalias2015}
likewise let the type exclude cases.

\begin{figure}
\small
\setlength{\jot}{2pt}
\begin{center}
\begin{tabular}{rcll}
\grammarGroup{Heads} \\
$h$ & $::=$ & $\ell$ & (Literal) \\
    & $\mid$ & $C$ & (Class) \\
$H$ & $::=$ & $\set{h_1, \ldots, h_n}$ & (Exclusion set)
   \\[3mm]
\grammarGroup{Shapes} \\
$k$ & $::=$ & $\shapeAny{\tau}{H}$ & (Rest) \\
    & $\mid$ & $\ell$ & (Literal) \\
    & $\mid$ & $\shapeAny{\patConstr{C}{k_1, \ldots, k_n}}{H}$ & (Constr) \\
    & $\mid$ & $\patTuple{k_1, \ldots, k_n}$ & (Tuple) \\
    & $\mid$ & $\shapeList{\tau}{k_1, \ldots, k_n}$ & (List) \\
    & $\mid$ & $\shapeDict{\tau}{\beta}{H}$ & (Dict) \\
$K$, $L$ & $::=$ & $\set{k_1, \ldots, k_n}$ & (Shape set)
\end{tabular}
\end{center}

\vspace{2mm}

\flushleft \shadebox{$\Sigma \vdash h : \tau$}
\begin{mathpar}
\small
\inferrule*[lab={\ruleName{head-literal}}]
{
  \Sigma \vdash \tyLiteral{\ell} \subty \tau
}
{
  \Sigma \vdash \ell : \tau
}
\and
\inferrule*[lab={\ruleName{head-class}}]
{
  \Sigma \vdash C \subty \tau
}
{
  \Sigma \vdash C : \tau
}
\and
\inferrule*[lab={\ruleName{head-length}}]
{
  n \text{ an integer}
}
{
  \Sigma \vdash n : \tyList{\tau}
}
\end{mathpar}

\vspace{2mm}

\flushleft \shadebox{$\Sigma \vdash k : \tau$}
\begin{mathpar}
\small
\inferrule*[lab={\ruleName{shape-rest}}]
{
  \Sigma \vdash h : \tau \quad (\forall h \in H) \\
  \tau \text{ neither union nor dictionary type}
}
{
  \Sigma \vdash \shapeAny{\tau}{H} : \tau
}
\and
\inferrule*[lab={\ruleName{shape-literal}}]
{
  \strut
}
{
  \Sigma \vdash \ell : \tyLiteral{\ell}
}
\and
\inferrule*[lab={\ruleName{shape-constr}}]
{
  \Sigma \vdash h : C \quad (\forall h \in H) \\
  \Sigma \vdash k_x : \sigma \quad (\forall \bind{x}{\sigma} \in \fields(\Sigma, C))
}
{
  \Sigma \vdash \shapeAny{\patConstr{C}{\seq{k}}}{H} : C
}
\and
\inferrule*[lab={\ruleName{shape-tuple}}]
{
  \Sigma \vdash k_i : \tau_i \quad (\forall i)
}
{
  \Sigma \vdash \patTuple{\seq{k}} : \tyTuple{\seq{\tau}}
}
\and
\inferrule*[lab={\ruleName{shape-list}}]
{
  \Sigma \vdash k_i : \tau \quad (\forall i)
}
{
  \Sigma \vdash \shapeList{\tau}{\seq{k}} : \tyList{\tau}
}
\and
\inferrule*[lab={\ruleName{shape-sub}}]
{
  \Sigma \vdash k : \sigma \\
  \Sigma \vdash \sigma \subty \tau
}
{
  \Sigma \vdash k : \tau
}
\and
\inferrule*[lab={\ruleName{shape-dict}}]
{
  \Sigma \vdash \beta(w) : \tau \quad (\forall w \in \dom(\beta)) \\
  \dom(\beta) \cap H = \varnothing
}
{
  \Sigma \vdash \shapeDict{\tau}{\beta}{H} : \tyDict{\tau}
}
\end{mathpar}
\caption{Shapes and shape typing}
\label{fig:shapes}
\end{figure}


A shape $k$ (\figref{shapes})
denotes a set of values of its type $\shapeTy(k)$ (\secref{types-auxiliary}). The form $\shapeAny{\tau}{H}$ denotes the
values of type $\tau$ whose outermost head is not among those in $H$. A head itself denotes a set of values: a
literal its value, a class its instances, and $\headList{n}$ the lists of length $n$. Head
typing $\Sigma \vdash h : \tau$ admits the heads whose values lie within $\tau$, and shape typing
$\Sigma \vdash k : \tau$ checks each excluded head and each component of a shape. The other forms are constructors, tuples, lists and dictionaries, each with shapes in place of
sub-patterns. A list shape also records its element type, which the element shapes do not determine. A
constructor shape carries an excluded set of its own: $\shapeAny{\patConstr{C}{\seq{k}}}{H}$ denotes the
instances of $C$, or of a subclass not below a class of $H$, whose fields, those of $C$, have the shapes
$\seq{k}$. The shapes of $\tau$ that remain once the heads in $H$ are excluded are
$\shapes(\Sigma, \tau, H)$: one shape per member of a union, and none at all when $H$ exhausts $\tau$. No shape carries a
union or dictionary type in the form $\shapeAny{\tau}{H}$, and $\shapes$ is the only source of that
form. A union has no head to exclude, so $\shapes$ splits it into its members, and a dictionary is
refined by its keys rather than by a head, so $\shapes$ gives it the dictionary form.

\begin{figure}
\small
\setlength{\jot}{2pt}
\flushleft \shadebox{$\Sigma; \Gamma \vdash k, p \matches K, L \binding \Delta$}
\begin{mathpar}
\small
\inferrule*[lab={\ruleName{pat-var}}]
{
  \strut
}
{
  \Sigma; \Gamma \vdash k, x \matches \set{k}, \varnothing \binding \set{\bind{x}{\shapeTy(k)}}
}
\and
\inferrule*[lab={\ruleName{pat-wild}}]
{
  \strut
}
{
  \Sigma; \Gamma \vdash k, \patWild \matches \set{k}, \varnothing \binding \varnothing
}
\and
\inferrule*[lab={\ruleName{pat-literal}}]
{
  \strut
}
{
  \Sigma; \Gamma \vdash \ell, \ell \matches \set{\ell}, \varnothing \binding \varnothing
}
\and
\inferrule*[lab={\ruleName{pat-as}}]
{
  \Sigma; \Gamma \vdash k, p \matches K, L \binding \Delta
}
{
  \Sigma; \Gamma \vdash k, p\;\pyinline{as}\;x \matches K, L \binding \Delta \disjunion \set{\bind{x}{\joinTy(\Sigma, \shapeTy(m) \mid m \in K)}}
}
\and
\inferrule*[lab={\ruleName{pat-split}}]
{
  \Sigma; \Gamma \vdash k, p \splits K', L' \\
  \Sigma; \Gamma \vdash K', p \matches K, L \binding \Delta
}
{
  \Sigma; \Gamma \vdash k, p \matches K, L \cup L' \binding \Delta
}
\and
\inferrule*[lab={\ruleName{pat-tuple}}]
{
  \Sigma; \Gamma \vdash \seq{k}, \seq{p} \matches K, L \binding \Delta
}
{
  \Sigma; \Gamma \vdash \patTuple{k_1, \ldots, k_n}, \patTuple{p_1, \ldots, p_n} \matches \set{\patTuple{\seq{m}} \mid \seq{m} \in K}, \set{\patTuple{\seq{k}\,'} \mid \seq{k}\,' \in L} \binding \Delta
}
\and
\inferrule*[lab={\ruleName{pat-list}}]
{
  \Sigma; \Gamma \vdash \seq{k}, \seq{p} \matches K, L \binding \Delta
}
{
  \Sigma; \Gamma \vdash \shapeList{\tau}{k_1, \ldots, k_n}, \patList{p_1, \ldots, p_n} \matches \set{\shapeList{\tau}{\seq{m}} \mid \seq{m} \in K}, \set{\shapeList{\tau}{\seq{k}\,'} \mid \seq{k}\,' \in L} \binding \Delta
}
\and
\inferrule*[lab={\ruleName{pat-dict}}]
{
  \seq{w} \text{ distinct} \\
  k_i = \beta(w_i) \quad (\forall i) \\
  \Sigma; \Gamma \vdash \seq{k}, \seq{p} \matches K, L \binding \Delta
}
{
  \Sigma; \Gamma \vdash \shapeDict{\tau}{\beta}{N}, \patDict{\bind{\seq{w}}{\seq{p}}} \matches \set{\shapeDict{\tau}{\beta \override \set{\bind{\seq{w}}{\seq{m}}}}{N} \mid \seq{m} \in K}, \set{\shapeDict{\tau}{\beta \override \set{\bind{\seq{w}}{\seq{k}\,'}}}{N} \mid \seq{k}\,' \in L} \binding \Delta
}
\and
\inferrule*[lab={\ruleName{pat-constr}}]
{
  \Gamma \vdash q : \Type{D} \\
  \Sigma \vdash C \subty D \\
  \fields(\Sigma, D) = \bind{\seq{x}}{\seq{\sigma}} \\
  \fieldMap(\Sigma, D, a) = \set{\bind{\seq{x}}{\seq{p}}} \\
  p'_i = \patWild \quad (\forall i) \\
  \Sigma; \Gamma \vdash \seq{k}, \seq{p} \mathbin{\concat} \seq{p}\,' \matches K, L \binding \Delta
}
{
  \Sigma; \Gamma \vdash \shapeAny{\patConstr{C}{\seq{k}}}{N}, \patConstr{q}{a} \matches \set{\shapeAny{\patConstr{C}{\seq{m}}}{N} \mid \seq{m} \in K}, \set{\shapeAny{\patConstr{C}{\seq{k}\,'}}{N} \mid \seq{k}\,' \in L} \binding \Delta
}
\end{mathpar}
\caption{Pattern $p$ matches shape $k$, giving matched shapes $K$, residual $L$ and bindings $\Delta$}
\label{fig:pat}
\end{figure}

\begin{figure}
\small
\setlength{\jot}{2pt}
\flushleft \shadebox{$\Sigma; \Gamma \vdash \seq{k}, \seq{p} \matches K, L \binding \Delta$}
\begin{mathpar}
\small
\inferrule*[lab={\ruleName{pat-seq}}]
{
  \Sigma; \Gamma \vdash k_i, p_i \matches K_i, L_i \binding \Delta_i \quad (\forall i) \\
  L = \textstyle\bigcup_i K_1 \times \cdots \times K_{i-1} \times L_i \times \set{k_{i+1}} \times \cdots \times \set{k_n}
}
{
  \Sigma; \Gamma \vdash \seq{k}, \seq{p} \matches K_1 \times \cdots \times K_n, L \binding \Delta_1 \disjunion \ldots \disjunion \Delta_n
}
\end{mathpar}

\vspace{2mm}

\flushleft \shadebox{$\Sigma; \Gamma \vdash L, p \matches K, L' \binding \Delta$}
\begin{mathpar}
\small
\inferrule*[lab={\ruleName{pat-shapes}}]
{
  \varnothing \neq M = \set{k \in L \mid \Sigma; \Gamma \vdash k, p \matches K_k, L_k \binding \Delta_k} \\
  \Delta = \set{\bind{x}{\joinTy(\Sigma, \Delta_k(x) \mid k \in M)} \mid x \in \binds(p)}
}
{
  \Sigma; \Gamma \vdash L, p \matches \textstyle\bigcup_{k \in M} K_k, \; \textstyle\bigcup_{k \in M} L_k \cup (L \setminus M) \binding \Delta
}
\end{mathpar}

\vspace{2mm}

\flushleft \shadebox{$\Sigma; \Gamma; \sigma \vdash \seq{p} \matches \seq{\Delta}, L$}
\begin{mathpar}
\small
\inferrule*[lab={\ruleName{pat-cases}}]
{
  L_0 = \shapesOf(\Sigma, \sigma, \varnothing) \\
  \Sigma; \Gamma \vdash p_i \agrees \sigma \quad (\forall i) \\
  \Sigma; \Gamma \vdash L_{i-1}, p_i \matches K_i, L_i \binding \Delta_i \quad (\forall i)
}
{
  \Sigma; \Gamma; \sigma \vdash \seq{p} \matches \seq{\Delta}, L_{n+1}
}
\end{mathpar}
\caption{Matching sequences, sets of shapes, and case patterns}
\label{fig:pat-seq}
\end{figure}


The matching judgement $\Sigma; \Gamma \vdash L, p \matches K, L' \binding \Delta$ divides an ordered set of shapes $L$ into
matched shapes $K$ and residual $L'$, binding $\Delta$ (\figref{pat}). It is defined from a
judgement on a single shape, which for a tuple, list, dictionary or constructor shape matches the sub-patterns
against the component shapes position by position. A pattern has no derivation against a shape it is unable to match, so
\ruleName{pat-shapes} takes the shapes matched by the pattern, of which there must be at least one, and passes the
unmatched shapes into the residual untouched. A variable pattern binds at $\shapeTy(k)$ for the shape $k$ it matches, and a
pattern named by \pyinline{as} at the join over the shapes it matched, so a case for one alternative of a
union leaves a later variable case bound at the remaining members alone. Bindings from the positions of a sequence, or from
the keys of a dictionary pattern, compose by disjoint union, so a pattern that binds a name twice has no
derivation.

\begin{figure}
\small
\setlength{\jot}{2pt}
\flushleft \shadebox{$\Sigma; \Gamma \vdash k, p \splits K, L$}
\begin{mathpar}
\small
\inferrule*[lab={\ruleName{split-literal}}]
{
  \ell \notin H \\
  \Sigma \vdash \tyLiteral{\ell} \subty \tau
}
{
  \Sigma; \Gamma \vdash \shapeAny{\tau}{H}, \ell \splits \set{\ell}, \shapesOf(\Sigma, \tau, H \cup \set{\ell})
}
\and
\inferrule*[lab={\ruleName{split-tuple}}]
{
  \strut
}
{
  \Sigma; \Gamma \vdash \shapeAny{\tyTuple{\tau_1, \ldots, \tau_n}}{\varnothing}, \patTuple{p_1, \ldots, p_n} \splits \set{\patTuple{\seq{k}} \mid \seq{k} \in \shapesOf(\Sigma, \seq{\tau})}, \varnothing
}
\and
\inferrule*[lab={\ruleName{split-list}}]
{
  n \notin H
}
{
  \Sigma; \Gamma \vdash \shapeAny{\tyList{\tau}}{H}, \patList{p_1, \ldots, p_n} \splits \set{\shapeList{\tau}{\seq{k}} \mid \seq{k} \in \shapesOf(\Sigma, \tau, \varnothing)^n}, \shapesOf(\Sigma, \tyList{\tau}, H \cup \set{n})
}
\and
\inferrule*[lab={\ruleName{split-key}}]
{
  w_i \notin \dom(\beta) \cup H \\
  w_j \in \dom(\beta) \quad (\forall j < i)
}
{
  \Sigma; \Gamma \vdash \shapeDict{\tau}{\beta}{H}, \patDict{\bind{w_1}{p_1}, \ldots, \bind{w_n}{p_n}} \splits \set{\shapeDict{\tau}{\beta \override \set{\bind{w_i}{k}}}{H} \mid k \in \shapesOf(\Sigma, \tau, \varnothing)}, \set{\shapeDict{\tau}{\beta}{(H \cup \set{w_i})}}
}
\and
\inferrule*[lab={\ruleName{split-class}}]
{
  \Gamma \vdash q : \Type{C} \\
  \ancestors(\Sigma, C) \cap H = \varnothing \\
  D = \meetTy{\Sigma}{\tau}{C} \\
  \fields(\Sigma, D) = \bind{\seq{x}}{\seq{\sigma}} \\
  H' = \set{h \in H \mid \Sigma \vdash h : D}
}
{
  \Sigma; \Gamma \vdash \shapeAny{\tau}{H}, \patConstr{q}{a} \splits \set{\shapeAny{\patConstr{D}{\seq{k}}}{H'} \mid \seq{k} \in \shapesOf(\Sigma, \seq{\sigma})}, \shapesOf(\Sigma, \tau, H \cup \set{C})
}
\and
\inferrule*[lab={\ruleName{split-subclass}}]
{
  \Gamma \vdash q : \Type{D} \\
  \Sigma \vdash D \subty C \\
  D \neq C \\
  \ancestors(\Sigma, D) \cap H = \varnothing \\
  \fields(\Sigma, D) = \fields(\Sigma, C) \mathbin{\concat} \bind{\seq{y}}{\seq{\sigma}} \\
  H' = \set{h \in H \mid \Sigma \vdash h : D}
}
{
  \Sigma; \Gamma \vdash \shapeAny{\patConstr{C}{\seq{k}}}{H}, \patConstr{q}{a} \splits \set{\shapeAny{\patConstr{D}{\seq{k} \mathbin{\concat} \seq{k}\,'}}{H'} \mid \seq{k}\,' \in \shapesOf(\Sigma, \seq{\sigma})}, \set{\shapeAny{\patConstr{C}{\seq{k}}}{(H \cup \set{D})}}
}
\end{mathpar}
\caption{Shape $k$ splits at the head tested by $p$, giving shapes $K$ and residual $L$}
\label{fig:split}
\end{figure}


A pattern matches a shape directly where the shape has the head tested by the pattern, and otherwise splits it first (\figref{split}): splitting takes the head
tested by the pattern, gives the shapes with that head, drawing the shapes of the component types from
$\shapes$, and leaves the shape with that head excluded as well. \ruleName{pat-split} matches the
pattern against the shapes the split gives and adds the shapes it leaves to the residual. Only the forms
with an excluded set split, and no head types at a tuple type, so $H$ is empty there and
\ruleName{split-tuple} leaves nothing.
A constructor pattern whose class lies below a
class of the excluded set matches nothing, because an earlier case for the superclass took every
instance.

A dictionary shape has a form of its own,
$\shapeDict{\tau}{\beta}{H}$, standing for dictionaries that hold each key of $\beta$ with a value of the
shape $\beta$ gives it, and hold no key of $H$. A dictionary pattern takes its keys in order, no key twice. A key in
neither $\beta$ nor $H$ splits the shape at that key (\ruleName{split-dict}), adding to the residual the
dictionaries that lack the key; once every key of the pattern is in $\beta$ the keys match as a sequence
(\ruleName{pat-dict}), each against the shape $\beta$ records for it. The empty dictionary pattern matches every
dictionary and leaves nothing.

\begin{lemma}[Preservation of shape type]
\label{lem:shape-typing}
If $\Sigma \vdash h : \tau$ for every $h \in H$ then $\Sigma \vdash k : \tau$ for each $k \in
\shapes(\Sigma, \tau, H)$. Moreover, if $\Sigma \vdash k : \tau$ for each $k \in L$ and $\Sigma; \Gamma \vdash L, p
\matches K, L' \binding \Delta$ then $\Sigma \vdash k : \tau$ for each $k \in K \disjunion L'$.
\end{lemma}

The case rules (\figref{typing-cases}) take the cases in the order $\dispatch$ tries them at run
time. A case that matches nothing is rejected, whether because the scrutinee can never take that shape or
because the earlier cases already cover it. The evaluation rules apply a tuple or list pattern only to a value
of its own kind, although Python would match either (\figref{eval-pat}); each case's pattern must
therefore be \emph{sequence-safe} at the scrutinee's type (\secref{types-auxiliary}), so a match on
$\tyList{\tyInt} \mathbin{|} \tyTuple{\tyInt, \tyInt}$ admits no tuple or list pattern, and a match on
$\tySized \mathbin{|} \tyTuple{\tyInt, \tyInt}$ no tuple pattern.

A union is exhausted
by covering each member, a class type by a pattern for the class itself, \pyinline{bool} by
\pyinline{True} and \pyinline{False}, and a dictionary type by the empty dictionary pattern, while \pyinline{int}, \pyinline{str} and lists have
unboundedly many heads and need a variable or wildcard. A literal-typed scrutinee is exhausted by covering
its literals, so \pyinline{x = 42} admits a match with a case for \pyinline{42} alone, while
\pyinline{x: int = 42} keeps the type open. Exhaustiveness need not come from a catch-all: over
$\tyTuple{\tyBool, \tyBool}$ the cases $\patTuple{\pyinline{True}, \patWild}$,
$\patTuple{\patWild, \pyinline{False}}$ and $\patTuple{\pyinline{False}, \pyinline{True}}$ leave nothing,
although each covers only part of the type.

\subsection{Modules and programs}
\label{sec:well-formedness-modules}

\emph{Static loading} types a module's body and computes its signature $\Gamma$, the judgement
$\Sigma \vdash_{\mathsf{load}} q : \Gamma \binding \Sigma'$ (\figref{typing-module}); it mirrors loading at run time
(\secref{operational-semantics}). The \ruleName{module} rule types the body as a top-level statement, and
the \ruleName{predefined} rule gives a predefined module the signature of
Definition~\ref{def:predefined}.
Both import forms statically load their target, together with its ancestors
(\figref{typing-import-prefix}). A plain import also records each of these modules as a member of its
parent (the static loads-as judgement), so that the binding created for the root of the target gives attribute
access along the whole chain. A from-import also loads any imported name that denotes a submodule. An
ancestor of the target is exempt from loading first if it is the importing module or one of its
ancestors, since at runtime it may still be loading. The exemption applies only to from-imports, because
for a plain import
the only ancestor that could never load first is the importing module itself, and an import naming a proper
descendant of the importing module is excluded outright\paperonly{ (\secref{no-load-dependent})}. Loading a module
may thus require other modules to load first; the rules are interpreted inductively, so a program with an
import cycle is ill-formed, because a module in a cycle could finish loading only after itself. Program
well-formedness (\figref{typing-module}) requires the main module to be well typed and hence,
recursively, every module the program can load; a module that is never imported is not checked.
The rule \ruleName{module} types a module body under the context of $\mathcal{P}(\pyinline{builtins})$,
overridden by the imported context, so the members of \pyinline{builtins} are in scope
unqualified.

\begin{figure}
\flushleft \shadebox{$\Sigma \vdash \Gamma, x \importsR \theta \binding \Sigma'$}
\begin{mathpar}
\small
\inferrule*[lab={\ruleName{imp-member}}]
{
  \Gamma \vdash x : \theta
}
{
  \Sigma \vdash \Gamma, x \importsR \theta \binding \Sigma
}
\and
\inferrule*[lab={\ruleName{imp-submodule}}]
{
  \bind{x}{\modStub{q}} \in \Gamma \\
  \Sigma \vdash_{\mathsf{load}} q : \Delta \binding \Sigma'
}
{
  \Sigma \vdash \Gamma, x \importsR \modLoaded{q}{\Delta} \binding \Sigma'
}
\end{mathpar}

\caption{Name $x$ imported from module signature $\Gamma$ has entry $\theta$}
\label{fig:typing-imports}
\end{figure}

\begin{figure}
\begin{subfigure}{\textwidth}
\flushleft \shadebox{$\Sigma; q \vdash \seq{\iota} : \Gamma \binding \Sigma'$}
\begin{mathpar}
\small
\inferrule*[lab={\ruleName{imports-nil}}]
{
  \strut
}
{
  \Sigma; q \vdash \varepsilon : \varnothing \binding \Sigma
}
\and
\inferrule*[lab={\ruleName{import}}]
{
  q' = x \cons \seq{x} \\
  q \not\properPrefixOf q' \\
  \Sigma \vdash_{\mathsf{load}} q' : \Delta \binding \Sigma_1 \\
  \Sigma_1 \vdash q', \modLoaded{q'}{\Delta} \loadsAs \theta \binding \Sigma_2 \\
  \Sigma_2; q \vdash \seq{\iota} : \Gamma \binding \Sigma'
}
{
  \Sigma; q \vdash \exImport{q'} \cons \seq{\iota} : \set{\bind{x}{\theta}} \extend \Gamma \binding \Sigma'
}
\and
\inferrule*[lab={\ruleName{from-import}}]
{
  \Sigma \vdash_{\mathsf{load}} q' : \Delta \binding \Sigma_0 \\
  q^\dagger_1 \properPrefixOf \cdots \properPrefixOf q^\dagger_m \text{ the } q^\dagger \properPrefixOf q' \text{ with } q^\dagger \not\prefixOf q \\
  \Sigma_{j-1} \vdash_{\mathsf{load}} q^\dagger_j : \Delta_j \binding \Sigma_j \quad (\forall j) \\
  \Sigma_{m+i-1} \vdash \Delta, x_i \importsR \theta_i \binding \Sigma_{m+i} \quad (\forall i) \\
  \Sigma_{m+n+1}; q \vdash \seq{\iota} : \Gamma \binding \Sigma'
}
{
  \Sigma; q \vdash \exFromImport{q'}{x_1, \ldots, x_{n+1}} \cons \seq{\iota} : \set{\bind{x_1}{\theta_1}, \ldots, \bind{x_{n+1}}{\theta_{n+1}}} \extend \Gamma \binding \Sigma'
}
\end{mathpar}

\subcaption{Import prefix $\seq{\iota}$ of module $q$ produces bindings $\Gamma$}
\end{subfigure}

\vspace{2mm}

\begin{subfigure}{\textwidth}
\flushleft \shadebox{$\Sigma \vdash q, \theta \loadsAs \theta' \binding \Sigma'$}
\begin{mathpar}
\small
\inferrule*[lab={\ruleName{loads-as-simple}}]
{
  \strut
}
{
  \Sigma \vdash x, \theta \loadsAs \theta \binding \Sigma
}
\and
\inferrule*[lab={\ruleName{loads-as-qualified}}]
{
  \Sigma \vdash_{\mathsf{load}} q : \Gamma' \binding \Sigma_1 \\
  \Sigma_1 \vdash q, \modLoaded{q}{\Gamma' \extend \set{\bind{x}{\theta}}} \loadsAs \theta' \binding \Sigma'
}
{
  \Sigma \vdash q.x, \theta \loadsAs \theta' \binding \Sigma'
}
\end{mathpar}
\subcaption{Reference $\theta$ for module $q$ loads as reference $\theta'$ for root of $q$, statically loading ancestors of $q$}
\end{subfigure}
\caption{Import typing}
\label{fig:typing-import-prefix}
\end{figure}

\begin{figure}
\flushleft \shadebox{$\Sigma \vdash_{\mathsf{load}} q : \Gamma \binding \Sigma'$}
\begin{mathpar}
\small
\inferrule*[lab={\ruleName{predefined}}]
{
  \mathcal{P}(q) = (\Gamma', \rho)
}
{
  \Sigma \vdash_{\mathsf{load}} q : \Gamma' \binding \Sigma
}
\and
\inferrule*[lab={\ruleName{module}}]
{
  \mathcal{M}(q) = \seq{\iota}\;t \\
  \Sigma; q \vdash \seq{\iota} : \Gamma \binding \Sigma_1 \\
  \mathcal{P}(\pyinline{builtins}) = (\Gamma', \rho) \\
  \Sigma_1; q; \Gamma' \mathbin{\override} \Gamma \vdash t : \Assigns{\Delta} \binding \Sigma' \\
  (\dom(\Gamma) \cup \dom(\Delta)) \cap \dom(\submods(q)) = \varnothing
}
{
  \Sigma \vdash_{\mathsf{load}} q : \submods(q) \override \Delta \binding \Sigma'
}
\end{mathpar}

\flushleft \shadebox{$\vdash \mathcal{M}$}
\begin{mathpar}
\small
\inferrule*[lab={\ruleName{program}}]
{
  \varnothing \vdash_{\mathsf{load}} \pymath{\_\_main\_\_} : \Gamma \binding \Sigma
}
{
  \vdash \mathcal{M}
}
\end{mathpar}
\caption{Module typing}
\label{fig:typing-module}
\end{figure}


The signature $\Gamma$ that static loading produces (\figref{typing-module}) comprises the classes
and variables a module defines, together with its intrinsic \pyinline{__name__} and a stub for each submodule,
never the names it imports: following PEP 484\footnote{\url{https://peps.python.org/pep-0484/}}, an imported
name is a private dependency, not re-exported as part of the interface. Submodule stubs and the module's own
bindings are disjoint, because the module rules exclude a module that binds a name that is also the name of
a submodule\paperonly{ (\secref{no-load-dependent})}.

